% Figura: problemas 4.10, 4.11 y 4.12. Condensadores planos con dieléctricos
% «en paralelo», «en serie» y con una lámina conductora intercalada.
\begin{figure}[t]
\centering
\begin{subfigure}[b]{0.3\textwidth}\centering
\begin{tikzpicture}[scale=0.8]
  \fill[dielecB] (0,1.7) rectangle (2.0,3.4);
  \fill[dielec] (0,0) rectangle (2.0,1.7);
  \draw[texto!70!fondo, line width=0.5pt] (0,1.7) -- (2.0,1.7);
  \draw[placa] (0,-0.1) -- (0,3.5);
  \draw[placa] (2.0,-0.1) -- (2.0,3.5);
  \draw[cargapos, line width=1.6pt] (-0.06,0) -- (-0.06,3.4);
  \node[etiq] at (1.0,3.05) {$\varepsilon_{r1}$};
  \node[etiq] at (1.0,1.35) {$\varepsilon_{r2}$};
  \draw[vecE] (0.45,2.35) -- (1.55,2.35);
  \draw[vecE] (0.45,0.65) -- (1.55,0.65);
  \node[etiq, anchor=east] at (-0.15,2.95) {$A_1$};
  \node[etiq, anchor=east] at (-0.15,2.3) {$Q_1$};
  \node[etiq, anchor=east] at (-0.15,1.25) {$A_2$};
  \node[etiq, anchor=east] at (-0.15,0.6) {$Q_2$};
  \draw[cota] (0,3.8) -- (2.0,3.8) node[midway, above, etiq] {$d$};
  % batería
  \draw[texto, line width=0.6pt] (0,-0.1) -- (0,-0.95) -- (0.8,-0.95);
  \draw[texto, line width=0.6pt] (1.2,-0.95) -- (2.0,-0.95) -- (2.0,-0.1);
  \draw[texto, line width=1.1pt] (0.85,-0.72) -- (0.85,-1.18);
  \draw[texto, line width=0.6pt] (1.12,-0.82) -- (1.12,-1.08);
  \node[etiq] at (1.0,-1.45) {$V_0$};
\end{tikzpicture}
\caption{Problema 4.10}
\end{subfigure}\hfill
\begin{subfigure}[b]{0.34\textwidth}\centering
\begin{tikzpicture}[scale=0.8]
  \fill[dielec] (0,0) rectangle (3.2,1.2);
  \draw[texto!70!fondo, line width=0.5pt] (0,1.2) -- (3.2,1.2);
  \draw[placa] (-0.1,0) -- (3.3,0);
  \draw[placa] (-0.1,2.6) -- (3.3,2.6);
  \node[etiq] at (2.95,0.6) {$\varepsilon$};
  \node[etiq] at (2.95,1.9) {$\eo$};
  \draw[vecD] (0.6,0.15) -- (0.6,1.05); \draw[vecD] (0.6,1.45) -- (0.6,2.35);
  \node[etiq, anchor=west, font=\footnotesize] at (0.68,1.9) {$\vD$};
  \draw[vecE] (1.6,0.4) -- (1.6,0.8); \draw[vecE] (1.6,1.35) -- (1.6,2.45);
  \node[etiq, anchor=west, font=\footnotesize] at (1.68,1.9) {$\vEext$};
  \node[etiqm, anchor=west, font=\footnotesize] at (1.68,0.6) {$\vEint$};
  \draw[cota] (-0.35,0) -- (-0.35,1.2) node[midway, left, etiq] {$a$};
  \draw[cota] (-0.35,1.2) -- (-0.35,2.6) node[midway, left, etiq] {$b$};
  % batería a la derecha
  \draw[texto, line width=0.6pt] (3.3,2.6) -- (3.95,2.6) -- (3.95,1.5);
  \draw[texto, line width=0.6pt] (3.95,1.1) -- (3.95,0) -- (3.3,0);
  \draw[texto, line width=1.1pt] (3.72,1.45) -- (4.18,1.45);
  \draw[texto, line width=0.6pt] (3.82,1.18) -- (4.08,1.18);
  \node[etiq, anchor=west] at (4.2,1.3) {$V_0$};
\end{tikzpicture}
\caption{Problema 4.11}
\end{subfigure}\hfill
\begin{subfigure}[b]{0.32\textwidth}\centering
\begin{tikzpicture}[scale=0.8]
  \fill[conductor] (0,1.2) rectangle (3.2,1.65);
  \draw[texto, line width=0.5pt] (0,1.2) rectangle (3.2,1.65);
  \draw[placa] (-0.1,0) -- (3.3,0);
  \draw[placa] (-0.1,3.0) -- (3.3,3.0);
  \node[etiq] at (0.45,0.6) {$\eo$};
  \node[etiq] at (0.45,2.33) {$\eo$};
  \node[etiq, anchor=west] at (3.35,3.0) {$+Q$};
  \node[etiq, anchor=west] at (3.35,1.73) {$-Q$};
  \node[etiq, anchor=west] at (3.35,1.12) {$+Q$};
  \node[etiq, anchor=west] at (3.35,0) {$-Q$};
  \draw[vecE] (1.2,2.85) -- (1.2,1.8); \draw[vecE] (1.2,1.05) -- (1.2,0.15);
  \node[etiq, anchor=west, font=\footnotesize] at (1.28,2.33) {$\sigma/\eo$};
  \node[etiq, anchor=west, font=\footnotesize] at (1.28,0.6) {$\sigma/\eo$};
  \node[font=\footnotesize, text=black] at (1.9,1.425) {$\vE=\vcero$};
  \draw[cota] (-0.35,0) -- (-0.35,1.2) node[midway, left, etiq] {$a$};
  \draw[cota] (-0.35,1.2) -- (-0.35,1.65) node[midway, left, etiq] {$e$};
  \draw[cota] (-0.35,1.65) -- (-0.35,3.0) node[midway, left, etiq] {$b$};
\end{tikzpicture}
\caption{Problema 4.12}
\end{subfigure}
\caption{Condensadores planos de área $A$ y separación $d$ parcialmente llenos. (a) Dos dieléctricos con la interfase
perpendicular a las placas: el campo es el mismo en ambos ($V_0/d$), pero la carga libre se reparte en
proporción a $\varepsilon_{ri}A_i$; el sistema equivale a dos condensadores en paralelo. (b) Interfase paralela a las
placas: $\vD$ es el mismo en ambas capas y el campo es menor en el dieléctrico; equivale a dos condensadores en
serie. (c) Una lámina conductora de espesor $e$ anula el campo en su interior y reduce la separación efectiva
a $d-e$, cualquiera que sea su posición.}
\label{fig:p410-412}
\end{figure}
